\section{Limitations and threats to validity}\label{sec:limits}

\paragraph{One model and GPU, two pools, two engines.} All registered SGLang evidence is for SGLang 0.5.22, Qwen2.5-7B-Instruct and one A100, at two pool sizes: the 300k-token pool limit and the 955,757-token pool SGLang sizes itself. The second answers the concern that the 300k-token pool limit created the cliff: with the native pool the cliff moved out and the registered test passed there too. vLLM 0.32.0 ran only with the 300k-token pool, on the same model and GPU. Its own closed curve and from-empty edge select the same window, $(1,\,2.0)\,\lref$, the same test points, 1.35 and 1.65 $\lref$, and the same cap of 20 as SGLang's 300k cell (\S\ref{sec:vllm}). The exploratory vLLM BURST runs preceded those measurements and had no NOBURST or CAPPED arm, so on their own they cannot exclude that a vLLM NOBURST run would also have congested at these rates within 8700\,s, or that a cap would not have recovered it. The registered vLLM paired test (\S\ref{sec:vllm}) uses the same seeds, reference rate, multiples and burst, so its BURST arm repeats the exploratory runs' workloads and does not replicate them independently. Larger models, parallelism, other GPUs and offloading are untested.

\paragraph{Synthetic sessions and clients.} Sessions are synthetic (random tokens, fixed tool outputs and output lengths, geometric turn counts, exponential 5\,s tools) and share no prefixes, which removes the case PEEK's scheduler targets. These parameters were fixed before we measured any trace, and they differ from the two trace sources of \S\ref{sec:robust}. Our mean tool time of 5\,s lies between theirs (1.40 and 8.51\,s), but their medians are 15--33 times shorter than ours, 63--85\% of their tool times are under 1\,s against 18\% of ours, and their coefficients of variation are 4.3--4.7 against 1. Their outputs per call are 2.7--5.9 times our 100 tokens, and TraceLab's mean prompt, 172k tokens, exceeds our 49,152-token context limit. Apart from one secondary mini-SWE-agent workload, the trace-driven runs of \S\ref{sec:robust} replace only the tool-time law, and only in the simulator. With no shared prefix a miss reuses nothing from other sessions. A shared 2000-token system prompt, cached once, would cut the prompt work of a session from 154,529 to 125,007 tokens when every turn misses and from about 9,600 to 7,600 when every follow-up hits, so the ratio of the two, about 16, would barely change. Production agent traces report a hit rate of about 55\% across turn boundaries and user idle gaps of minutes~\citep{liu2026copilot}, while 93\% of our turns hit in the good state (not the same metric). The simulator suggests that real tool-time laws keep the trap but move the window and the cap; the engine has not tested this. Our clients never retry or abandon a request, and the 7200\,s request timeout never fired, although trapped turns took 163--237\,s on average from send to completion at 300k and 321--511\,s with the native pool. Retries would add load in the trapped state and abandonment would remove it, so either could change whether the engine recovers; neither was tested.

\paragraph{Few runs, finite horizons, nominal rates, uncontrolled factors.} The closed curves and the from-empty sweeps have one run per point, so the peak's location (20 against 24 at 300k and on vLLM, 64 against 80 native) is within noise and single runs near the cliff land on either side of it. The registered tests, the grid and the vLLM runs share the same three seeds, so their runs are not independent across rates, pools or engines (\S\ref{sec:harness}). All rates in this paper are nominal. Each from-empty sweep used a single arrival seed, whose realized rates were 6--13\% above nominal, and in the evaluation windows of the pre-registered SGLang tests the runs labelled 1.35 and 1.65 $\lref$ received 1.27--1.45 and 1.61--1.75 $\lref$ (native, 2.5 and 4.0: 2.48--2.73 and 3.89--4.28; vLLM BURST: 1.24--1.45 and 1.56--1.76). ``Does not leave'' means within 7500\,s on the engine and 30,000\,s in the simulator; the 3000\,s horizon of the from-empty runs cannot see slow escapes, so the true edges may be below 2.0 $\lref$ (both engines) and 5.5 $\lref$ (native). The native cap rests on that coarse curve: $\mu(80)$ was within 0.35\% of $\mu(64)$, and the cliff lies somewhere between 80 and 96 with no run in between. The registered smallest-within-5\% rule gives 64 whichever of 64 and 80 is the true peak (the registration still flags the cap as ambiguous), and with no grid point between 48 and 64 it could not have found a smaller cap; a rule that takes the argmax could as well have picked 80, which may sit within one grid step of the cliff; how far a cap of 64 is from the cliff was not measured. The registered rule loses margin at both ends of the window: the slope test near $\lmiss$, and the token test at high multiples, where the trapped engine's prefill rate rises above the all-miss rate (1.106--1.124 times it at 4.0 $\lref$ with the native pool, against a tolerance of 1.15), so at somewhat higher rates the rule would stop calling a stuck run trapped. In both SGLang closed curves the cliff coincides with a change of job and node (at 300k between $n=32$ and 48, native between 80 and 96; Table~\ref{tab:stages}), so neither curve excludes a node or job effect on where the cliff falls. The collapsed side ran on different nodes in the two curves, which argues against one slow node; on vLLM the hit fraction crossed 0.5 between two runs of one job on one node ($n=28$ and 32); and both registered SGLang tests, which ran on all three nodes with consistent outcomes, argue against a node effect on the trap. The SGLang server seed was not pinned, and the load client runs on the server's node.

\paragraph{Model and simulators.} Eq.~\eqref{eq:fluid} is used only for the window, the recovery threshold and the capped drain time, which the engine met within $-4\%$ to $+52\%$ given the population present at the burst, and it is not quantitative in congestion (\S\ref{sec:fluid}). Proposition~\ref{prop:A} is a limit statement about the closed-population throughput $\mu(n)$ as $n$ grows and says nothing about any finite population. It covers request-level FCFS and, on our reading of the engines' sources, vLLM's FCFS with its preemption and SGLang's stock LPM, which stops re-sorting its queue when more than 128 requests wait; it does not cover program-level FCFS, pure resident-first orders or eviction rules that foresee tool times. It assumes that long-run rates exist, private prefixes, at most one outstanding request per session, no offloading of evicted KV, i.i.d.\ tool times with finite mean that are independent of everything else and unknown to the engine in advance, a bounded number of running requests, a bounded admission rate, and an engine that in congestion computes prompt tokens no faster than with caching disabled; it gives equality only if the engine computes them as fast. With zero hits our engine computed them at 1.03 times the all-miss rate near the stationary turn mix and at 0.99--1.12 times it over the windows we measured, so (A3) and (A4) hold on it only to within a few percent at the stationary mix and to within 12\% over those windows. The proof and Lemma~\ref{lem:waits} were checked by AI review agents and have not been refereed. The open system's recovery threshold rests on a heuristic work balance and, without a monotone tail of $\mu$, can lie below $\lmiss$; that recovery from deep congestion without a cap needs a rate below $\lmiss$ is untested on the engine; pure LPM's congested rate is a conjecture; and the capped system's stochastic stability is not proved. The calibrated simulator failed its level test at $n=48$, reserves whole turns, omits partial prefixes and the refresh, frees a finished session's KV at once (an oracle: a real engine keeps it as evictable cache), and understated LPM. In its primary mode it put the native from-empty edge at least one grid step too high in nominal multiples, against a single real run whose arrivals exceeded the real peak, and for vLLM it predicted a later, gentler cliff than the engine showed (hit fraction 0.46--0.53 against 0.08 at $n=32$). The mechanism comparison uses an older simulator whose assumed costs are about 1.5 times the fitted ones, so its absolute rates and recovery thresholds do not carry over to the engine. Its pin rows are fixed, configured TTLs, not Continuum's cost model; its PEEK rows are ablations and a projection of the released scheduler; and its CONCUR row is a surrogate of the paper's control equation.
